\documentclass[11pt]{amsart}

\usepackage[T1]{fontenc}
\usepackage{lmodern}
\usepackage{microtype}
\usepackage{amsmath,amssymb,amsthm}
\usepackage[hidelinks]{hyperref}

\hypersetup{
  pdftitle={A Proof of a Conjecture of Balanescu and Cimpoeas on the Hilbert
Depth of Intersections of Monomial Primes}
}

\newtheorem{theorem}{Theorem}
\newtheorem{lemma}[theorem]{Lemma}

\newcommand{\hdepth}{\operatorname{hdepth}}
\newcommand{\NN}{\mathbb{Z}_{\ge 0}}

\title[Hilbert depth of intersections of monomial primes]
{A Proof of a Conjecture of B\u{a}l\u{a}nescu and Cimpoea\c{s} on the
Hilbert Depth of Intersections of Monomial Primes}
\date{}

\subjclass[2020]{Primary 13C15; Secondary 13F20, 05A18}
\keywords{Hilbert depth, monomial prime ideal, Hilbert series,
palindromic polynomial}

\begin{document}

\begin{abstract}
Let $S=K[x_1,\ldots,x_N]$, and let $I$ be the intersection of monomial
prime ideals generated by pairwise disjoint blocks of variables of sizes
$n_1,\ldots,n_r$, where $N=n_1+\cdots+n_r$. B\u{a}l\u{a}nescu and
Cimpoea\c{s} proved
\[
\hdepth(I)\le \left\lfloor\frac{N+r}{2}\right\rfloor
\]
and conjectured equality. We prove the matching lower bound, thereby
resolving their Conjecture~3.4. The proof combines Uliczka's criterion
with a coefficientwise-positivity lemma for centered shifts of
palindromic polynomials.
\end{abstract}

\maketitle

\section{Introduction}

Let $K$ be a field. Partition the variables of
$S=K[x_1,\ldots,x_N]$ into nonempty blocks $B_1,\ldots,B_r$ of sizes
$n_1,\ldots,n_r$, and set
\[
P_i=(x_j:j\in B_i),
\qquad
I=P_1\cap\cdots\cap P_r,
\qquad
N=n_1+\cdots+n_r.
\]
B\u{a}l\u{a}nescu and Cimpoea\c{s} proved the upper bound
\[
\hdepth(I)\le \left\lfloor\frac{N+r}{2}\right\rfloor
\]
and obtained equality for $r=2$ and whenever at most one of the $n_i$ is
even \cite[Theorems~2.9 and~3.3]{BC}. They proposed equality in general
as Conjecture~3.4. Our contribution is the reverse inequality for
arbitrary positive block sizes.

\begin{theorem}\label{thm:main}
With the notation above,
\[
\hdepth(I)=\left\lfloor\frac{N+r}{2}\right\rfloor.
\]
In particular, Conjecture~3.4 of B\u{a}l\u{a}nescu and Cimpoea\c{s} is
true.
\end{theorem}

\section{Proof}

A formal power series is called nonnegative if all its coefficients are
nonnegative.

\begin{lemma}\label{lem:palindrome}
Let
\[
Q(u)=\sum_{j=0}^{D}q_j u^j,
\qquad q_j=q_{D-j}\ge 0,
\]
and put $\delta=\lceil D/2\rceil$. Then
\[
(1-t)^{-\delta}Q(1-t)
\]
is a nonnegative formal power series.
\end{lemma}

\begin{proof}
We first note that, for integers $b\ge a\ge0$,
\begin{equation}\label{eq:atom}
(1-t)^a+(1-t)^{-b}
\end{equation}
is nonnegative. The case $b=0$ is immediate. If $b\ge1$, then for every
$m\ge0$ its coefficient of $t^m$ is
\[
(-1)^m\binom{a}{m}+\binom{b+m-1}{m}.
\]
This is nonnegative when $m$ is even or $m>a$. If $m$ is odd and
$1\le m\le a$, then $b+m-1\ge a$, so
\[
\binom{b+m-1}{m}\ge\binom{a}{m}.
\]
Thus \eqref{eq:atom} is nonnegative.

If $D=2\delta$, palindromicity gives
\[
u^{-\delta}Q(u)
=q_{\delta}+\sum_{e=1}^{\delta}q_{\delta+e}
  (u^e+u^{-e}).
\]
If $D=2\delta-1$, it gives
\[
u^{-\delta}Q(u)
=\sum_{e=0}^{\delta-1}q_{\delta+e}
  (u^e+u^{-e-1}).
\]
Substituting $u=1-t$ and applying \eqref{eq:atom} with
$(a,b)=(e,e)$ in the first case and $(a,b)=(e,e+1)$ in the second proves
the claim.
\end{proof}

For a finitely generated graded $S$-module $M$, write
\[
H_M(t)=\sum_{k\ge0}\dim_K(M_k)t^k.
\]
Uliczka's criterion states that
\begin{equation}\label{eq:uliczka}
\hdepth(M)
=\max\bigl\{d\in\NN:(1-t)^dH_M(t)
\text{ is nonnegative}\bigr\}
\end{equation}
\cite[Theorem~3.2]{Uliczka}. B\u{a}l\u{a}nescu and Cimpoea\c{s} invoke this
criterion in the squarefree reformulation \cite[Theorem~2.4]{BCK}, which is
derived from it; we use it directly.

\begin{proof}[Proof of Theorem~\ref{thm:main}]
A monomial lies in $I$ precisely when its support meets every block.
Consequently,
\begin{align*}
H_I(t)
&=\prod_{i=1}^{r}\left((1-t)^{-n_i}-1\right)\\
&=\frac{t^rQ(1-t)}{(1-t)^N},
\end{align*}
where
\[
Q(u)=\prod_{i=1}^{r}[n_i]_u,
\qquad
[n]_u=1+u+\cdots+u^{n-1}.
\]
The polynomial $Q$ has nonnegative coefficients, is palindromic, and has
degree
\[
D=\sum_{i=1}^{r}(n_i-1)=N-r.
\]
Set
\[
q=\left\lfloor\frac{N+r}{2}\right\rfloor,
\qquad
\delta=N-q=\left\lceil\frac{D}{2}\right\rceil.
\]
Then
\[
(1-t)^qH_I(t)
=t^r(1-t)^{-\delta}Q(1-t),
\]
which is nonnegative by Lemma~\ref{lem:palindrome}. Hence
\eqref{eq:uliczka} yields $\hdepth(I)\ge q$.

The reverse inequality was proved in \cite[Theorem~3.3]{BC}. For
completeness, it also follows from the first coefficient after the
initial term. If $d>q$, then
\begin{align*}
[t^{r+1}](1-t)^dH_I(t)
&=(N-d)Q(1)-Q'(1)\\
&=Q(1)\left(\frac{N+r}{2}-d\right)<0.
\end{align*}
Indeed, palindromicity of degree $D$ gives
$Q'(1)=\frac{D}{2}Q(1)$, and $Q(1)=\prod_i n_i>0$. Thus no $d>q$ is
admissible in \eqref{eq:uliczka}, and $\hdepth(I)=q$.
\end{proof}

\medskip
\noindent\textbf{Status and scope.}
Theorem~\ref{thm:main} establishes the equality proposed as Conjecture~3.4 of
\cite{BC} for an arbitrary number of blocks and arbitrary positive block
sizes. Only the lower bound is proved here; the upper bound is due to
B\u{a}l\u{a}nescu and Cimpoea\c{s} \cite[Theorem~3.3]{BC}, and the
alternative derivation given above is included for completeness rather than
as a new result. The equality is recorded as a conjecture in the published
version of \cite{BC} and is still stated as a conjecture in the most recent
arXiv revision of that paper (v6, 14 June 2026); we have found no proof of
the general case in the literature, but our search is not exhaustive, and the
statement is offered subject to external prior-art review. Nothing here bears
on the separate conjecture of \cite{BC} approximating $\hdepth(S/I)$, which is
untouched. The Stanley depth of $I$ is likewise not determined here: the
argument concerns Hilbert depth only, and we make no claim about
$\operatorname{sdepth}(I)$.

\begin{thebibliography}{9}

\bibitem{BC}
S.~B\u{a}l\u{a}nescu and M.~Cimpoea\c{s},
\emph{On the Hilbert depth of certain monomial ideals and applications},
U.P.B. Sci. Bull., Series A \textbf{87} (2025), no.~4, 53--66;
\href{https://arxiv.org/abs/2306.11015}{arXiv:2306.11015}.

\bibitem{BCK}
S.~B\u{a}l\u{a}nescu, M.~Cimpoea\c{s} and C.~Krattenthaler,
\emph{On the Hilbert depth of monomial ideals},
Comm. Algebra (2026), 1--19, published online 20 January 2026;
\href{https://doi.org/10.1080/00927872.2025.2608666}
{doi:10.1080/\allowbreak 00927872.2025.2608666};
\href{https://arxiv.org/abs/2306.09450}{arXiv:2306.09450}.

\bibitem{Uliczka}
J.~Uliczka,
\emph{Remarks on Hilbert series of graded modules over polynomial rings},
Manuscripta Math. \textbf{132} (2010), 159--168;
\href{https://doi.org/10.1007/s00229-010-0341-9}
{doi:10.1007/s00229-010-0341-9}.

\end{thebibliography}

\end{document}